\documentclass[12pt]{article}

\usepackage{amsmath,amssymb,hyperref}
\topmargin -.5cm
\textheight 21cm
\oddsidemargin -.125cm
\textwidth 17cm

\begin{document}

\begin{center}
{\Large\bf Local TeX Transcription for arXiv:2504.12290}

\vspace{0.3cm}
{\large Diffeomorphism in Closed String Field Theory}
\end{center}

\noindent
This file is a local TeX transcription of the portions of
\href{https://arxiv.org/abs/2504.12290}{arXiv:2504.12290}
used in the SFT notes.  The sandbox in this environment cannot resolve
\texttt{arxiv.org}, so the original source archive could not be downloaded directly into
\texttt{/references}.  The text below records the key statements from the ar5iv rendering that are
used in the note
\texttt{notes/axion\_omega\_quantization/axion\_omega\_quantization.tex}.

\section*{Bibliographic data}

M.~Cho, S.~Collier, J.~Sandor and X.~Yin,
``Diffeomorphism in Closed String Field Theory,''
\href{https://arxiv.org/abs/2504.12290}{arXiv:2504.12290}.

\section*{Key statements from Sections 4.1--4.4}

\paragraph{Riemann-normal-coordinate family.}
The paper introduces, for each point on the target manifold, a Riemann normal coordinate system and a
corresponding family of string fields related by diffeomorphism gauge transformations.  The basic point is
that one should keep track simultaneously of the solution in many local coordinate patches rather than
around a single origin.

\paragraph{Extended BV system.}
The equations of motion for the background string field and the gauge-transport equations that move
between nearby Riemann normal coordinates are combined into a single extended BV system.  In the text this
is described as an enlarged string field valued in differential forms on the target manifold together with
an extended BRST operator.  The resulting equation simultaneously encodes:
\begin{itemize}
\item the ordinary SFT equation of motion,
\item the covariance of the family of fields under infinitesimal diffeomorphisms,
\item and the integrability condition for transporting between different coordinate origins.
\end{itemize}

\paragraph{Flux background.}
For the compact flux example analyzed in Section 4.3, the paper emphasizes that one can choose a global
orthonormal frame adapted to the isometries, so that the relevant background fields become independent of
the base point on the target manifold.  The ghost-number-one string fields then realize the global
isometry algebra.

\paragraph{Spectrum of fluctuations (Section 4.4).}
The fluctuation problem is formulated so that the fluctuation string fields satisfy both an on-shell
condition and a covariance condition under the global isometries.  The fluctuations are organized into
isometry representations.  Applying the representation condition twice shows that the target-space
Laplacian agrees with the quadratic Casimir, to the relevant order in the weak-curvature expansion.
Matching to the target-space Laplacian spectrum then implies that the allowed representations appear in the
BRST cohomology with the expected multiplicities.

\paragraph{Crucial conceptual statement.}
The paper stresses that this spectral result relies crucially on the fact that the extended fields know
about the \emph{entire target space}, not only the local patch visible in an expansion in $x/R$.  This is
presented as the main advantage of the extended BV system for recovering global information such as KK
quantization.

\section*{Conclusion statements used in the notes}

The conclusion makes two points that are directly relevant for the current AdS application:
\begin{itemize}
\item the perturbative SFT description around a single coordinate origin only sees a local patch and does
      not detect the full topology or global isometries of the target space;
\item the extended BV system remedies this by transporting the solution along the manifold and by making
      the fluctuation problem globally covariant with respect to the isometries.
\end{itemize}
The paper explicitly says that quantization was not visible in the straightforward perturbative analysis,
and that understanding how the extended formalism imposes the expected restrictions on the spectrum is one
of its main uses.

\end{document}
